\section{Computational Methods}
\label{ch:methods}

\subsection{Fixation-gated circular delayed-response task}

Each circular-task trial comprised pre-cue fixation, a circular cue, a silent
memory delay, and a response period (Fig.~\ref{fig:tasks}A). Stimulus angle
$\theta\in[0,2\pi)$ was encoded over 32 evenly spaced preferred angles
$\mu_j$ as
\begin{equation}
 c_j(\theta)=\exp\{\kappa[\cos(\theta-\mu_j)-1]\},\qquad \kappa=8.
\end{equation}
The fixation input and output remained high until the response cue. Circular
outputs were deliberately silent before response and then reported the
remembered angle. Pre-cue and cue durations varied during training; evaluation
used clean delays of 10, 20, 40, and 80 model steps. A separate delay-20
condition inserted a five-step irrelevant circular cue midway through the
delay. Training sampled one clean and one distractor-present homogeneous batch
exactly once per shuffled two-batch block. Distractor filtering was therefore a
learned property of the evaluated networks rather than an out-of-distribution
challenge.

\subsection{Context-cued N-back task}

The N-back task presented 20 sequential items drawn from six stimulus
identities (Fig.~\ref{fig:tasks}B). Two context channels instructed either
0-back or 2-back behaviour. In 0-back, the network classified whether the
current item matched a fixed target identity. In 2-back, it classified whether
the current item matched the item presented two positions previously.
Sequences contained six matches; the 2-back generator additionally controlled
one-back lures. The first two items were excluded from scored outcomes. The
output consisted of match and non-match logits.

\begin{figure*}[t]
  \centering
  \includegraphics[width=0.98\textwidth]{task_schematic.png}
  \caption{Complementary working-memory tasks.
  \textbf{A}, The fixation-gated circular task required a continuous cue angle
  to be maintained across variable delays and reported after the go cue;
  balanced clean/distractor training made resistance to an irrelevant
  delay-period cue a learned behaviour.
  \textbf{B}, Context inputs switched the same N-back network between a
  fixed-target 0-back rule and a 2-back updating rule. The tasks supplied
  settling, delay, distractor, and updating-load contrasts rather than a single
  performance level.}
  \label{fig:tasks}
\end{figure*}

\subsection{Network architecture}

Both task families used the same rate-style continuous-time recurrent network
(CTRNN) family: a leaky recurrent hidden layer with a linear readout. The
model is not spiking, not Dale-constrained, and not intended as a
biophysically detailed cortical circuit. It is a controlled dynamical testbed
with an explicit time constant so that maintenance stability can be varied
computationally.

Let $x_t\in\mathbb{R}^{N_{\mathrm{in}}}$ be the input at discrete time $t$,
$h_t\in\mathbb{R}^{N_h}$ the hidden state, and $y_t\in\mathbb{R}^{N_{\mathrm{out}}}$
the readout. The native recurrence mixes afferent and recurrent drive,
\begin{equation}
 z_t = W_{\mathrm{in}}x_t+b_{\mathrm{in}}
      +W_{\mathrm{rec}}h_{t-1}+b_{\mathrm{rec}},
 \label{eq:drive}
\end{equation}
then applies a leaky update with nonlinearity after mixing,
\begin{equation}
 h_t = \tanh\!\left[(1-\alpha)h_{t-1}+\alpha z_t\right],
 \qquad \alpha=\frac{dt}{\tau}.
 \label{eq:native}
\end{equation}
Here $dt=20$ and $\tau=100$, so $\alpha=0.2$. The factor $(1-\alpha)$ is the
carried-state weight and $\alpha$ the drive weight. Because $\tanh$ acts after
leak mixing rather than only on $z_t$, later gain and persistence operators are
translations of this implemented update, not exact maps of a Song-style
$\tau\dot h=-h+f(z)$ form. Hidden states were initialised at zero at the start
of each trial. Outputs were a linear readout of the hidden state,
\begin{equation}
 y_t = W_{\mathrm{out}}h_t+b_{\mathrm{out}}.
 \label{eq:readout}
\end{equation}

Unless otherwise stated, $N_h=64$ and the hidden nonlinearity was $\tanh$.
The circular model used $N_{\mathrm{in}}=N_{\mathrm{out}}=33$ (32 tuned
channels plus fixation). The N-back model used eight inputs (six stimulus
identities and two rule-context channels) and two outputs (non-match and match
logits). Learned weights were frozen before all perturbation analyses.

\subsection{Training and checkpoint replication}

Throughout, a \emph{checkpoint} means one independently trained network with
frozen weights, identified by its random seed. Checkpoint, not trial or
technical replicate, is the replication unit: summary statistics are formed
across trained seeds, and within-seed trial averages are treated as
measurement precision rather than independent scientific replicates.

Circular networks were trained for 4,000 Adam updates (learning rate
$10^{-3}$, batch size 64) with Yang-style weighted fixation and response loss
on a balanced clean/distractor mixture; recurrent training noise standard
deviation was $0.05$. Six circular seeds were trained; five met preconfigured
competence limits (clean mean error $\le10^\circ$, distractor mean error
$\le15^\circ$, fixation accuracy $\ge0.90$) and were retained (seeds 20260731,
20260732, 20260733, 20260735, and 20260736). Seed 20260734 exceeded those
limits (clean $15.75^\circ$, distractor $17.46^\circ$) and was excluded before
the perturbation screen. The ten N-back checkpoints
(seeds 20260912--20260921) were drawn from a competence-screened pool trained
with staged 0-back then 2-back curriculum under Adam (learning rate $10^{-3}$,
no recurrent training noise). Final N-back retention required accuracy
$\ge0.95$, discriminability $\ge0.90$, and 2-back lure accuracy $\ge0.90$;
stage-1 0-back acquisition used stricter gates (accuracy $\ge0.98$,
discriminability $\ge0.95$).

These gates are operational competence thresholds, not fitted human effect
sizes. Chance circular error is of order $90^\circ$, so clean $\le10^\circ$
requires unambiguous angle memory while remaining softer than the few-degree
errors typical of retained seeds; distractor $\le15^\circ$ allows a modest
filtering cost without accepting near-failure. Fixation $\ge0.90$ keeps the
network in a reliable go regime so latency contrasts are not confounded by
fixation collapse. N-back gates demand near-ceiling fixed-target performance,
clearly above-chance updating, and lure resistance, so load contrasts reflect
selective updating cost rather than a weak baseline.

\subsection{Prior baseline and structured-noise background}

Earlier project stages established that continuous circular delayed response is
learnable in the same CTRNN family and that noise structure matters for
representational robustness. A categorical tanh delay baseline reached perfect
held-out accuracy across tested delays, while a tuned circular baseline
supported stable angle decoding with delay-period slowing of hidden-state
speed. Those reports remain methodological background rather than the
psilocybin-signature screen.

A five-seed structured-noise study on Yang-style fixation-gated circular
checkpoints compared matched-RMS independent Gaussian, temporal AR(1), and
topology-correlated AR(1) hidden-state noise. At equal energy, impairment
followed independent $<$ temporal $<$ topology-correlated structure, grew with
delay, accumulated during maintenance, and coincided with collapse of the
circular memory geometry. That hierarchy is retained as a robustness result
and as motivation for treating unstructured Gaussian noise as the generic
control in Section~\ref{ch:results-specificity}. Temporally persistent and
topology-aligned noise are not used as primary psychedelic candidates in the
present signature screen: they destroy representations efficiently, but they
are not the literature-preferred mechanism class for selective behavioural
dissociations.

\subsection{Candidate perturbations}

We evaluated seven single-operator families against each checkpoint's native
baseline (Table~\ref{tab:perturbations}). Operators are competing computational
interventions motivated by, but not equated with, psychedelic or working-memory
mechanisms. Because the native update applies $\tanh$ after mixing drive with
carried state, no operator is an exact biological response-gain map; each is an
explicit RNN translation. Unless noted, every non-neutral operator was applied
at every timestep of the full frozen forward pass for the entire trial or
sequence; learned weights were never updated. The only timed exception is P3b,
which is confined to the irrelevant circular-cue window. Write
$z_t^{\mathrm{w}}=W_{\mathrm{in}}x_t+W_{\mathrm{rec}}h_{t-1}$ for the
weight-derived drive. Symbols are shared across perturbations:
$x_t$ is input, $h_t$ hidden state, and $z_t$ pre-activation drive at time
$t$; $W_{\mathrm{in}}$ and $W_{\mathrm{rec}}$ are input and recurrent weights;
$b_{\mathrm{in}}$ and $b_{\mathrm{rec}}$ are biases; and $\alpha=dt/\tau$ is
the native leak factor. In P3a/P3b, $x_t^{\mathrm{stim}}$ and
$x_t^{\mathrm{ctrl}}$ denote stimulus versus fixation/rule-control channels,
with matching input submatrices $W_{\mathrm{in}}^{\mathrm{stim}}$ and
$W_{\mathrm{in}}^{\mathrm{ctrl}}$. Scalar $g$ denotes uniform gain, vector
$\mathbf{g}$ fixed unit-wise gain (P2), $\varepsilon_t\sim\mathcal{N}(0,I)$
Gaussian noise scaled by $\sigma$ (P5), $p$ persistence scale (P6), and $s$
timescale factor with derived $\alpha'$ (P7). Unless noted, learned biases
stay outside multiplicative gains.

\subsubsection{P1 Synaptic-drive gain}
Herzog et al.\ operationalised 5-HT2A effects as response gain on synaptic
current inside a nonlinearity \cite{herzog2023}; REBUS likewise links altered
precision to postsynaptic gain \cite{carhartharris2019}. In this CTRNN the
closest available analogue multiplies weight-derived drive while leaving
biases unscaled:
\begin{equation}
 h_t=\tanh\!\left[
 (1-\alpha)h_{t-1}+\alpha\bigl(g z_t^{\mathrm{w}}+b_{\mathrm{in}}+b_{\mathrm{rec}}\bigr)
 \right].
 \label{eq:p1}
\end{equation}
This is synaptic-drive gain, not an exact $F$--$I$ slope change and not a pure
carried-state manipulation.

\subsubsection{P2 Heterogeneous drive gain}
Receptor-density dependence in Herzog et al.\ implies spatially nonuniform gain
rather than one global scalar \cite{herzog2023}. We therefore replaced $g$ by a
fixed positive unit-wise vector $\mathbf{g}$ drawn from a log-normal
distribution and normalised to population mean one:
\begin{equation}
 h_t=\tanh\!\left[
 (1-\alpha)h_{t-1}+\alpha\bigl(\mathbf{g}\odot z_t^{\mathrm{w}}+b_{\mathrm{in}}+b_{\mathrm{rec}}\bigr)
 \right].
 \label{eq:p2}
\end{equation}
Three fixed vectors were treated as technical replicates and averaged within
checkpoint; they are not independent trained seeds.

\subsubsection{P3a Sensory-input gain}
REBUS predicts increased influence of afferent or bottom-up signals when
high-level precision is relaxed \cite{carhartharris2019}. As a task-level
translation, only stimulus-channel input weights were scaled, leaving fixation
or rule-context channels unchanged:
\begin{equation}
\begin{aligned}
 z_t^{\mathrm{in}} &=
 g\,W_{\mathrm{in}}^{\mathrm{stim}}x_t^{\mathrm{stim}}
 +W_{\mathrm{in}}^{\mathrm{ctrl}}x_t^{\mathrm{ctrl}}
 +b_{\mathrm{in}},\\
 h_t &= \tanh\bigl[(1-\alpha)h_{t-1}
 +\alpha(z_t^{\mathrm{in}}+W_{\mathrm{rec}}h_{t-1}+b_{\mathrm{rec}})\bigr].
\end{aligned}
 \label{eq:p3a}
\end{equation}
This tests sensory sensitivity, not receptor-mapped response gain.

\subsubsection{P3b Distractor-window gain}
Human psilocybin work shows particular vulnerability when competing information
must be ignored \cite{carter2005}. P3b applies Eq.~\eqref{eq:p3a} only during
the irrelevant circular cue and is therefore a circular-task manipulation
check rather than a complete cross-task mechanism candidate.

\subsubsection{P4 Recurrent gain}
Recurrent excitation is a canonical substrate of delay-period maintenance
\cite{wang2013nmda}. To separate recurrent weighting from sensory drive we
scaled only the recurrent weight contribution:
\begin{equation}
\begin{aligned}
 z_t^{(4)} &=
 W_{\mathrm{in}}x_t+b_{\mathrm{in}}
 +g W_{\mathrm{rec}}h_{t-1}+b_{\mathrm{rec}},\\
 h_t &= \tanh\bigl[(1-\alpha)h_{t-1}+\alpha z_t^{(4)}\bigr].
\end{aligned}
 \label{eq:p4}
\end{equation}
Recurrent bias remained outside the gain. P4 is a mechanistic contrast within
the trained RNN, not a direct REBUS or 5-HT2A claim.

\subsubsection{P5 Gaussian state noise}
Independent Gaussian noise is retained in the project catalogue as a generic
disruption control against which selective signatures must later be shown to
exceed nonspecific damage. Noise was added to the pre-activation,
\begin{equation}
 h_t=\tanh\!\left[
 (1-\alpha)h_{t-1}+\alpha z_t+\sigma\varepsilon_t
 \right],
 \quad \varepsilon_t\sim\mathcal{N}(0,I),
 \label{eq:p5}
\end{equation}
in the exploratory pilot, but P5 was outside the candidate-only screen reported
in Section~\ref{ch:results-screen}. Section~\ref{ch:results-specificity} is
reserved for the matched-cost comparison.

\subsubsection{P6 State persistence}
Trained working-memory RNNs can store information in slow or persistent latent
dynamics \cite{ghazizadeh2021,wang2013nmda}. P6 asks whether weakening
retention of the carried state, without rescaling incoming drive, can produce
selective behavioural costs. The update is asymmetric:
\begin{equation}
 h_t=\tanh\!\left[p(1-\alpha)h_{t-1}+\alpha z_t\right],
 \label{eq:persistence}
\end{equation}
where $p$ is the persistence scale. This is an operational RNN intervention,
not a pre-specified biological mapping from psilocybin to reduced persistence;
$p=0.95$ entered the present screen as the exploratory pilot nominee.

\subsubsection{P7 Conserved effective time constant}
To test whether P6 effects are merely those of any timescale change, P7
rescales the leaky-integrator coefficients while preserving their sum. With
timescale factor $s$,
\begin{equation}
 \alpha'=\min(\alpha/s,1),\qquad
 h_t=\tanh\!\left[(1-\alpha')h_{t-1}+\alpha'z_t\right].
 \label{eq:p7}
\end{equation}
Thus $s>1$ slows the effective timescale and $s<1$ speeds it. Within the
evaluated grid, $\alpha/s<1$, so the minimum was inactive. P7 is a
mechanistic control for conserved timescale change, not a psychedelic
parameter claim.

\begin{table*}[t]
\caption{Candidate perturbations and fixed evaluation grids.
Neutral settings reproduce the native forward pass. Distractor-window gain was
defined only for the circular distractor condition; all other operators except
P5 were evaluated in both task families. P5 is listed for catalogue continuity
but was not part of the candidate-only screen.}
\label{tab:perturbations}
\centering
\begin{tabular}{lll}
\toprule
ID & Computational intervention & Strengths \\
\midrule
P1 & Scalar gain on input and recurrent weight drive
& 0.90--1.10 \\
P2 & Fixed mean-one unit-wise log-normal gain
& log-SD 0.00--0.30 \\
P3a & Gain on stimulus-channel input contribution
& 0.80--1.20 \\
P3b & P3a restricted to irrelevant cue window
& 1.00--1.50 \\
P4 & Scalar gain on recurrent weight contribution
& 0.90--1.10 \\
P5 & Additive pre-activation noise (catalogue)
& specificity chapter \\
P6 & Scale $p$ on carried-state coefficient
& 0.90--1.10 \\
P7 & Conserved update with $\alpha'=\min(\alpha/s,1)$
& 0.80--1.25 \\
\bottomrule
\end{tabular}
\end{table*}

\subsection{Evaluation and behavioural signature}

Each circular checkpoint, operator, strength, condition, and delay used 1,024
trials. Each N-back checkpoint, operator, strength, and load used 1,024
20-item sequences. Native and perturbed evaluations reused identical generated
task samples within checkpoint and cell.

Circular response accuracy was mean absolute angular error. Post-cue settling
was the first response step at which angular error was below $15^\circ$ and
population-vector amplitude exceeded 50\% of the native reference amplitude.
Trials that did not settle were capped at the 25-step response duration,
yielding restricted mean settling time. Settling was interpreted only when
fixation accuracy was at least 0.90 and at least 80\% of trials settled.
The $15^\circ$ settling threshold is a coarse latency criterion distinct from
continuous accuracy; half-native amplitude excludes collapsed bumps; and the
80\% settled floor marks latency invalid when most trials never settle.
``Slowing with preservation'' required a positive settling change at clean
delay 20 while proportional angular-error impairment remained at or below
0.20, encoding selective slowing without large accuracy loss.

For circular condition $c$, proportional impairment was
\begin{equation}
 I_c=\frac{E_{\mathrm{perturbed},c}-E_{\mathrm{native},c}}
 {E_{\mathrm{native},c}},
\end{equation}
where $E$ is mean angular error. Delay selectivity was
$I_{\mathrm{delay\,80}}-I_{\mathrm{delay\,10}}$; distractor selectivity was
$I_{\mathrm{distractor\,20}}-I_{\mathrm{clean\,20}}$.

N-back discriminability was $D=H-F$, where $H$ and $F$ are hit and false-alarm
rates. Condition-normalised impairment was
\begin{equation}
 J_c=\frac{D_{\mathrm{native},c}-D_{\mathrm{perturbed},c}}
 {D_{\mathrm{native},c}},
\end{equation}
and load selectivity was $J_{\mathrm{2back}}-J_{\mathrm{0back}}$. Positive
values indicate disproportionate high-load impairment.

The candidate-screen analysis was descriptive. For each operator-strength
setting we reported the checkpoint mean, SD, and a student $t$ 95\% interval
across seeds, together with the fraction of seeds in the predicted direction.
A setting met the complete descriptive pattern when the mean
slowing-with-preservation, distractor, and N-back load contrasts were positive
and at least a simple majority of checkpoints agreed on each (three of five
circular; six of ten N-back); delay selectivity was supporting evidence only.
That majority rule generalises the exploratory pilot's $\ge$2-of-3 criterion
and is implemented in the summary code; it was not restated as a numeric
threshold in the 1,024-trial preregistration text. Cross-operator comparisons
are uncontrolled for overall perturbation magnitude, so a near-miss may fail
because of dose as well as mechanism. No participant-level inference,
confirmatory $p$ values, or multiple-comparison claims were made for the
screen.
